\origpage{83}{N. Robertson, P. Seymour / Journal of Combinatorial Theory, Series B 89 (2003) 77--108}
\textbf{Proof.} Let us suppose that $(A, B)$ is not titanic. Hence we may choose subhypergraphs $X_1, X_2, X_3$ of $A$ such that $X_1 \cup X_2 \cup X_3 = A$ and $E(X_1), E(X_2), E(X_3)$ are mutually disjoint, and
\begin{align*}
|(V_1 \cup V_2) \cap V_3| &< |W_1 \cup W_2| \\
|(V_2 \cup V_3) \cap V_1| &< |W_2 \cup W_3| \\
|(V_3 \cup V_1) \cap V_2| &< |W_3 \cup W_1|
\end{align*}
where $V(X_i) = V_i$ and $V(X_i \cap B) = W_i$ $(i = 1, 2, 3)$. Suppose that $(X_1, X_2 \cup X_3 \cup B) \notin \mathscr{T}$. Then either $(X_2 \cup X_3 \cup B, X_1) \in \mathscr{T}$ or $(X_1, X_2 \cup X_3 \cup B)$ has order $\ge \theta$; and in either case, since $(B, A)$ is linked to $\mathscr{T}$, we deduce that $(X_1, X_2 \cup X_3 \cup B)$ has order at least that of $(B, A)$. Hence
\begin{equation*}
|V_1 \cap (V_2 \cup V_3 \cup V(B))| \ge |V(A \cap B)|,
\end{equation*}
that is,
\begin{equation*}
|V_1 \cap (V_2 \cup V_3)| + |W_1 \setminus (W_2 \cup W_3)| \ge |W_2 \cup W_3| + |W_1 \setminus (W_2 \cup W_3)|
\end{equation*}
contrary to our assumption. Hence $(X_1, X_2 \cup X_3 \cup B) \in \mathscr{T}$ and similarly $(X_2, X_3 \cup X_1 \cup B)$, $(X_3, X_1 \cup X_2 \cup B) \in \mathscr{T}$. It follows that $(X_1 \cup B, X_2 \cup X_3) \notin \mathscr{T}$ by (T2), since $(X_1 \cup B) \cup X_2 \cup X_3 = G$; and $(X_2 \cup X_3, X_1 \cup B) \notin \mathscr{T}$ since $(X_2 \cup X_3) \cup X_1 \cup B = G$; and so $(X_1 \cup B, X_2 \cup X_3)$ has order $\ge \theta$; that is,
\begin{align*}
\theta &\le |(V_1 \cup V(B)) \cap (V_2 \cup V_3)| = |(V_2 \cup V_3) \cap V_1| + |(W_2 \cup W_3) \setminus W_1| \\
&< |W_2 \cup W_3| + |(W_2 \cup W_3) \setminus W_1| = 2|(W_2 \cup W_3) \setminus W_1| + |(W_2 \cup W_3) \cap W_1|.
\end{align*}
By summing this and the two similar inequalities, we obtain
\begin{align*}
3\theta &< 2|(W_2 \cup W_3) \setminus W_1| + 2|(W_3 \cup W_1) \setminus W_2| + 2|(W_1 \cup W_2) \setminus W_3| \\
&\quad + |(W_2 \cup W_3) \cap W_1| + |(W_3 \cup W_1) \cap W_2| + |(W_1 \cup W_2) \cap W_3| \\
&= 4|W_1 \cup W_2 \cup W_3| - |W_1 \cap W_2 \cap W_3| \\
&\le 4|V(A \cap B)|.
\end{align*}
Hence the order of $(A, B)$ is $> 3\theta/4$, a contradiction, and so our initial assumption that $(A, B)$ is not titanic was false. This completes the proof of 3.3. $\square$

Let $\mathscr{L}$ be a location in $G$. The \textit{order} of $\mathscr{L}$ is the maximum order of the members of $\mathscr{L}$ (or $0$ if $\mathscr{L} = \emptyset$). For each $(A, B) \in \mathscr{L}$ let $e(A, B)$ be a new element, and let $H$ be the hypergraph with
\begin{equation*}
V(H) = V(M(G, \mathscr{L})),
\end{equation*}
\begin{equation*}
E(H) = E(M(G, \mathscr{L})) \cup \{e(A, B) \colon (A, B) \in \mathscr{L}\},
\end{equation*}
where for $e \in E(M(\mathscr{L}))$ its ends are as in $G$, and for $(A, B) \in \mathscr{L}$ the ends of $e(A, B)$ are the elements of $V(A \cap B)$. This is a hypergraph by 2.3.2, and we call it the \textit{heart} of $\mathscr{L}$. We define the \textit{tree-width} of $\mathscr{L}$ to be the tree-width of $H$.
